% from ega4-9-fr.tex
\oldpage[IV-3]{60}
ouverts affines partout denses, et les nombres $\dim((U_j)_\eta)$ sont tous
égaux à $n=\dim(X_\eta)$ (4.1.1.3)~; on peut par suite se borner au cas où $X$
est affine, donc aussi $X_\eta$. Il existe alors, en vertu de (4.1.2), un
ouvert non vide $W$ de $X$ tel que $W\supset X_\eta$, et un
$\mathbf{k}(\eta)$-morphisme fini surjectif
$h:W_\eta\to\mathbf V_{\mathbf{k}(\eta)}^n
(=\operatorname{Spec}(\mathbf{k}(\eta)[T_1,\ldots,T_n]))$~; en appliquant
(8.8.2, (i)) et la méthode de (8.1.2, \emph{a)}), on en déduit (en remplaçant
au besoin $S$ par un voisinage de $\eta$) que l'on a $h=g_\eta$, où
$g:W\to\mathbf V_A^n(=\operatorname{Spec}(A[T_1,\ldots,T_n]))$ est un
$S$-morphisme, et on peut supposer ce morphisme fini et surjectif en vertu de
(8.10.5, (vi) et (x)). On en conclut que pour tout $s\in S$, le morphisme
$g_s:W_s\to\mathbf V_{\mathbf{k}(s)}^n$ est fini et surjectif, donc que
$\dim(W_s)=n$ (4.1.2).

\subsection{Propriétés constructibles de morphismes de préschémas algébriques}

\begin{proposition}[9.3.1]
\label{IV.9.3.1-fr}
Soit $\mathbf P(X,k)$ une propriété constructible (resp. ind-constructible)
de préschémas algébriques. Désignons par $\mathbf P'(f,X,Y,k)$ la relation
suivante~: $f:X\to Y$ est un $k$-morphisme de $k$-préschémas algébriques tel
que pour tout $y\in Y$, on ait la propriété
$\mathbf P(f^{-1}(y),k(y))$. Alors $\mathbf P'$ est une propriété
constructible (resp. ind-constructible).
\end{proposition}

En effet, comme $\mathbf P$ vérifie la condition (9.2.1, (i)), il en est de
même de $\mathbf P'$ en vertu de la transitivité des fibres (I, 3.6.4)~; par
ailleurs le fait que $\mathbf P'$ vérifie la condition (9.2.1, (ii)) résulte
de (9.2.4), en raison de la remarque (9.2.5).

\begin{proposition}[9.3.2]
\label{IV.9.3.2-fr}
Désignons par $\mathbf P$ une des propriétés suivantes d'un $k$-morphisme
$f:X\to Y$ de $k$-préschémas algébriques~:
\begin{enumerate}[label=(\roman*)]
  \item $f$ est surjectif.
  \item $f$ est quasi-fini.
  \item $f$ est radiciel.
  \item Pour tout $y\in Y$, $\dim(f^{-1}(y))$ appartient à $\Phi$ (notation
  de (9.2.6)).
\end{enumerate}
Alors $\mathbf P$ est constructible.
\end{proposition}

Cela résulte aussitôt de (9.3.1) et (9.2.6) si l'on tient compte de ce que $f$
est de type fini (1.5.4, (v)), de la caractérisation des morphismes radiciels
(I, 3.5.8), et de celle des morphismes quasi-finis (II, 6.2.2).

\begin{proposition}[9.3.3]
\label{IV.9.3.3-fr}
Supposons vérifiées les hypothèses de (8.8.1) dont nous gardons les
notations~; supposons en outre que $S_\alpha$ soit quasi-compact, $X_\alpha$
et $Y_\alpha$ de présentation finie sur $S_\alpha$, et soit
$f_\alpha:X_\alpha\to Y_\alpha$ un $S_\alpha$-morphisme. Soit $\mathbf P$ une
propriété ind-constructible de morphismes de préschémas algébriques. Pour tout
$s\in S$ (resp. $s_\lambda\in S_\lambda$) posons
$X_s=X\times_S\operatorname{Spec}(k(s))$,
$Y_s=Y\times_S\operatorname{Spec}(k(s))$,
$f_s=f\times1:X_s\to Y_s$ (resp.
$X_{\lambda,s_\lambda}=X_\lambda\times_{S_\lambda}
\operatorname{Spec}(k(s_\lambda))$,
$Y_{\lambda,s_\lambda}=Y_\lambda\times_{S_\lambda}
\operatorname{Spec}(k(s_\lambda))$,
$f_{\lambda,s_\lambda}=f_\lambda\times1:
X_{\lambda,s_\lambda}\to Y_{\lambda,s_\lambda}$). Alors, afin que pour
tout $s\in S$ on ait la propriété $\mathbf P(f_s)$, il faut et il suffit
qu'il existe $\lambda\geqslant\alpha$ tel que pour tout
$s_\lambda\in S_\lambda$, on ait $\mathbf P(f_{\lambda,s_\lambda})$.
\end{proposition}

En effet, soit $E$ (resp. $E_\lambda$) l'ensemble des $s\in S$ (resp.
$s_\lambda\in S_\lambda$) tels que la propriété $\mathbf P(f_s)$ (resp.
$\mathbf P(f_{\lambda,s_\lambda})$) soit vraie~; il résulte de
(9.2.2, (iv)) que l'on a
$E_\mu=u_{\lambda\mu}^{-1}(E_\lambda)$ pour $\lambda\leqslant\mu$, et
$E=u_\lambda^{-1}(E_\lambda)$~; en outre, en vertu de (9.2.3), $E$ (resp.
$E_\lambda$) est ind-constructible dans $S$ (resp. $S_\lambda$)~; la
proposition résulte donc de (8.3.4) appliqué à la partie ind-constructible
$E$ de $S$.

